\documentclass[10pt, compress]{beamer}

\usetheme{glasgow}

\usepackage{booktabs}
\usepackage[scale=2]{ccicons}
\usepackage{minted}
\usepackage{bookmark}
\usepgfplotslibrary{dateplot}

\usemintedstyle{trac}

% ($ (A)!r!(B) $) the location of images to be used
\graphicspath{{src/}}

%% Customisation
% \newcommand{\V}[1]{\v} % vectors \v{c}
% \renewcommand{\v}[1]{\mathbf{#1}} % vectors
\newcommand{\ti}[1]{\tilde{#1}} % spectral representation
\newcommand{\tnsr}[1]{\underline{\underline{#1}}}

% Symbols
\renewcommand{\O}{\omega}  % omega
\newcommand{\E}{\varepsilon}  % epsilon
\renewcommand{\u}{\mu}  % mu
\newcommand{\p}{\rho}  % rho
\newcommand{\x}{\times}  % times
\renewcommand{\inf}{\infty}  % infinity
\newcommand{\infint}{\int\limits_{-\inf}^\inf} % integral by R
\newcommand{\e}{\mathrm{e}} % Straight-up exponential
\renewcommand{\j}{{j}\mkern1mu} % Straight-up exponential
\newcommand{\iu}{\mathrm{i}\mkern1mu}

\newcommand\ddfrac[2]{\frac{\displaystyle #1}{\displaystyle #2}}

\title{High Frequency Communication Systems}
\subtitle{Lecture 4}
\date{Spring 2021}
\author{Hasan T Abbas \& Qammer H Abbasi}
% \institute{}

\begin{document}

\maketitle

%%%%%%%%%%%%%%%%%%%%%%%%%%%%%%%%%%%%%%%%%%
%%%%%%%%%%%%%%%%%%%%%%%%%%%%%%%%%%%%%%%%%%
%%%%%%%%%%%%%%%%%%%%%%%%%%%%%%%%%%%%%%%%%%
\begin{frame}[fragile]
  \frametitle{Lecture Outline}
\begin{outline}[itemize]
  \1 Transmission Lines Theory \& Analysis
  \1 The Smith Chart 
  \1 Load Mismatching
  \1 The Magic of Quarter-wave Transformer
\end{outline}
\end{frame}
%%%%%%%%%%%%%%%%%%%%%%%%%%%%%%%%%%%%%%%%%%
%%%%%%%%%%%%%%%%%%%%%%%%%%%%%%%%%%%%%%%%%%
%%%%%%%%%%%%%%%%%%%%%%%%%%%%%%%%%%%%%%%%%%
\begin{frame}[fragile]
\frametitle{Transmission Line Theory}
\begin{outline}
  \1 A technique to bridge the gap between basic circuit analysis and electromagnetic fields theory
   \2 Lots of similaritites and analogies
    \1 Commonly used to design microwave devices and circuits
\end{outline}
\end{frame}

\begin{frame}
  \frametitle{TL \& EM Theory - Duality}
  \begin{columns}[T] % align columns
  \begin{column}{.3\textwidth}
\begin{outline}
  \1 In wave scattering problems, we place the co-ordinate system at the boundary
  \1 In transmission lines, we use the two different lines as the boundary
  \1 For open problems, the load is considered to be at infinity.
  \end{outline}
   \end{column}
 \begin{column}[T]{.7\textwidth}
    \begin{figure}
      \centering
          \includegraphics[width=.9\textwidth]{circuit.pdf}
      \caption{The analogy between wave problems and transmission lines.}
    \end{figure}
      \end{column}%
\end{columns}
  \end{frame}

\end{document}
